% !TEX program = xelatex
% 공통수학2 · Ⅲ. 함수와 그래프 · 1. 함수 · 01 함수 · 2/2차시
% 학생용 활동지 (답안 없음)
\documentclass[11pt,a4paper]{article}
\usepackage{kotex}
\usepackage{iftex}
\ifPDFTeX\else
  \setmainhangulfont{Noto Serif CJK KR}
  \setsanshangulfont{Noto Sans CJK KR}
\fi
\usepackage[margin=2.1cm]{geometry}
\usepackage{parskip}
\usepackage{enumitem}
\usepackage{array}
\usepackage{tabularx}
\usepackage{xcolor}
\usepackage{amssymb}
\usepackage{tikz}
\usepackage[most]{tcolorbox}
\usepackage{fancyhdr}
\usepackage{titlesec}

\definecolor{goldc}{HTML}{B07C22}
\definecolor{goldstrong}{HTML}{8A5F16}
\definecolor{indigoc}{HTML}{2B3B57}
\definecolor{indigosoft}{HTML}{6E82A6}
\definecolor{linec}{HTML}{D6DACB}
\definecolor{surface2}{HTML}{E4E8DC}
\definecolor{writeline}{HTML}{C7CCBA}
\definecolor{inksoft}{HTML}{52604F}

\titleformat{\section}{\normalfont\sffamily\bfseries\large\color{indigoc}}{\thesection}{0.6em}{}
\titlespacing{\section}{0pt}{14pt}{6pt}
\newtcolorbox{goalbox}{enhanced, breakable, colback=white, colframe=goldc, boxrule=0.5pt, leftrule=3pt, arc=2pt, left=10pt, right=10pt, top=8pt, bottom=8pt}
\newtcolorbox{sheetbox}{enhanced, breakable, colback=white, colframe=linec, boxrule=0.6pt, arc=3pt, left=11pt, right=11pt, top=9pt, bottom=9pt}
\newcommand{\writespace}[1]{%
  \par\nobreak\vspace{3pt}
  \foreach \n in {1,...,#1} {%
    \noindent\textcolor{writeline}{\rule{\linewidth}{0.4pt}}\par\vspace{0.62cm}%
  }
}
\newcommand{\fillblank}[1]{\makebox[#1]{\hrulefill}}
\pagestyle{fancy}\fancyhf{}\renewcommand{\headrulewidth}{0pt}
\fancyfoot[L]{\footnotesize\color{inksoft} 공통수학2 · 2026학년도 2학기 · 지평선고등학교}
\fancyfoot[R]{\footnotesize\color{inksoft} 다음 시간 --- 02 합성함수}

\newcommand{\mapthreefour}[1]{%
  \begin{tikzpicture}[scale=0.58, baseline]
    \draw[thick, indigosoft] (0,0) ellipse (0.9 and 1.9);
    \draw[thick, indigosoft] (2.8,0) ellipse (0.9 and 1.9);
    \node[circle, fill=indigoc, inner sep=1.4pt, label=left:{\footnotesize $1$}] (x1) at (0,1.2) {};
    \node[circle, fill=indigoc, inner sep=1.4pt, label=left:{\footnotesize $2$}] (x2) at (0,0) {};
    \node[circle, fill=indigoc, inner sep=1.4pt, label=left:{\footnotesize $3$}] (x3) at (0,-1.2) {};
    \node[circle, fill=goldstrong, inner sep=1.4pt, label=right:{\footnotesize $1$}] (y1) at (2.8,1.8) {};
    \node[circle, fill=goldstrong, inner sep=1.4pt, label=right:{\footnotesize $3$}] (y2) at (2.8,0.6) {};
    \node[circle, fill=goldstrong, inner sep=1.4pt, label=right:{\footnotesize $5$}] (y3) at (2.8,-0.6) {};
    \node[circle, fill=goldstrong, inner sep=1.4pt, label=right:{\footnotesize $7$}] (y4) at (2.8,-1.8) {};
    \node[above] at (0,2.1) {\footnotesize $X$};
    \node[above] at (2.8,2.1) {\footnotesize $Y$};
    #1
  \end{tikzpicture}}

\newcommand{\gridaxes}{%
  \draw[step=1cm, linec, very thin] (-2.2,-2.2) grid (2.2,2.2);
  \draw[->, thick, indigosoft] (-2.2,0) -- (2.5,0) node[right, font=\small]{$x$};
  \draw[->, thick, indigosoft] (0,-2.2) -- (0,2.5) node[above, font=\small]{$y$};}

\begin{document}

\noindent
\begin{minipage}[b]{0.62\linewidth}
  {\sffamily\footnotesize\color{indigosoft} 공통수학2 · Ⅲ. 함수와 그래프 · 1. 함수}\\[2pt]
  {\sffamily\bfseries\Large 01 함수 \textperiodcentered\ 34차시}
\end{minipage}\hfill
\begin{minipage}[b]{0.36\linewidth}
  \raggedleft\small\color{inksoft}
  반\ \fillblank{1.6cm} \quad 번호\ \fillblank{1.6cm}\\[4pt]
  이름\ \fillblank{3.6cm}
\end{minipage}

\vspace{4pt}
\noindent{\color{goldc}\rule{\linewidth}{2.2pt}}
\vspace{10pt}

\begin{goalbox}
\textbf{오늘의 목표}\\
함수의 그래프를 판별하고, 일대일함수·일대일대응·항등함수·상수함수를 구별할 수 있다.
\vspace{4pt}
\small\color{inksoft}
$\square$ 자신 있음 \quad $\square$ 조금 헷갈림 \quad $\square$ 처음 봄 \hfill (수업 전 나의 상태에 표시)
\end{goalbox}

\section{생각 열기 --- 그림을 좌표로}

\begin{sheetbox}
지난 시간에 배운 함수 $f:X\to Y$는 순서쌍 $(x,f(x))$ 전체의 집합, 즉 좌표평면 위의 점들의 모임으로도 나타낼 수 있다. $X=\{1,2,3\}$, $f(x)=2x-1$일 때, 세 점 $(1,f(1))$, $(2,f(2))$, $(3,f(3))$을 좌표로 각각 써 보자.
\writespace{2}

이 점들을 그래프의 `함수다움'과 연결 지어, 만약 어떤 $x$값에 점이 두 개 이상 있다면 무엇이 문제일지 한 문장으로 써 보자.
\writespace{2}
\end{sheetbox}

\section{개념 정리 --- 함수의 그래프와 판정법}

\begin{sheetbox}
함수 $f:X\to Y$에서 순서쌍 $(x,f(x))$ 전체의 집합을 함수 $f$의 \fillblank{2.4cm}라 한다.

\vspace{6pt}
\textbf{수직선 판정법} --- 좌표평면의 어떤 도형이 함수의 그래프가 되려면, $x$축에 \fillblank{2.4cm}인 직선(수직선)을 그었을 때 그 도형과 \fillblank{2.6cm}에서만 만나야 한다.

\vspace{8pt}
\noindent
\begin{tikzpicture}[scale=0.5]
  \gridaxes
  \draw[thick, goldc, domain=-1.85:1.85, smooth, variable=\t] plot({\t},{0.5*\t*\t-1});
  \draw[dashed, inksoft] (1,-2.2) -- (1,2.2);
  \node[below] at (0,-2.6) {\footnotesize ① $y=\frac12x^2-1$};
\end{tikzpicture}\hfill
\begin{tikzpicture}[scale=0.5]
  \gridaxes
  \draw[thick, goldc] (0,0) circle (1.5);
  \draw[dashed, inksoft] (0.8,-2.2) -- (0.8,2.2);
  \node[below] at (0,-2.6) {\footnotesize ② $x^2+y^2=2.25$};
\end{tikzpicture}

\vspace{4pt}
①, ② 중 함수의 그래프인 것은 \fillblank{1.6cm}이다. 그 이유를 써 보자.
\writespace{2}
\end{sheetbox}

\section{개념 정리 --- 여러 가지 함수}

\begin{sheetbox}
함수 $f:X\to Y$에 대하여

\begin{enumerate}[leftmargin=1.6em, itemsep=4pt]
  \item \textbf{일대일함수} --- $X$의 서로 다른 원소에 대응하는 함숫값이 \fillblank{2.6cm}, 즉 $x_1\neq x_2$이면 \fillblank{2.6cm}인 함수.
  \item \textbf{일대일대응} --- 일대일함수이면서 \fillblank{2cm}과 \fillblank{2cm}이 같은 함수.
  \item \textbf{항등함수} --- $X$의 각 원소 $x$에 그 자신 $x$를 대응시키는 함수. 즉 $f(x)=$ \fillblank{1.4cm}
  \item \textbf{상수함수} --- $X$의 모든 원소에 $Y$의 단 하나의 원소 $c$만을 대응시키는 함수. 즉 $f(x)=$ \fillblank{1.4cm}
\end{enumerate}
\end{sheetbox}

\section{연습 문제}

\begin{sheetbox}
\textbf{1.} $X=\{1,2,3\}$, $Y=\{1,3,5,7\}$, $f(x)=2x-1$의 대응을 아래 그림에 화살표로 나타내고, 이 함수가 일대일함수인지, 일대일대응인지 판별하시오.

\vspace{4pt}
\centering
\mapthreefour{}
\writespace{3}

\raggedright\vspace{4pt}
\textbf{2.} $X=\{1,2,3\}$일 때, 항등함수 $f(x)=x$와 상수함수 $g(x)=2$의 함숫값을 각각 구하시오.
\writespace{3}

\vspace{4pt}
\textbf{3.} 다음 중 일대일함수인 것을 모두 고르고, 아닌 것은 반례를 하나씩 제시하시오.\\
\hspace*{1.6em}⑴ $f(x)=3x+1$ \quad ⑵ $f(x)=x^2$ \quad ⑶ $f(x)=-2x$ \quad ⑷ $f(x)=5$
\writespace{4}

\vspace{4pt}
\textbf{4. (수평선 판정법)} $y=\frac12x^2-1$의 그래프는 일대일함수의 그래프인가? 임의의 수평선을 그어 확인하고 이유를 쓰시오.
\writespace{3}
\end{sheetbox}

\section{사고의 가시화 --- See · Think · Wonder}

\noindent{\footnotesize\color{inksoft} 오늘 배운 `여러 가지 함수'를 떠올리며 아래 세 칸을 채워 보자.}\\[6pt]
\begin{tabularx}{\linewidth}{@{}XXX@{}}
\textbf{\footnotesize\color{indigoc} See --- 무엇을 보았나?} &
\textbf{\footnotesize\color{indigoc} Think --- 무엇을 생각했나?} &
\textbf{\footnotesize\color{indigoc} Wonder --- 무엇이 궁금한가?} \\[4pt]
\end{tabularx}
\noindent
\begin{minipage}[t]{0.32\linewidth}\writespace{3}\end{minipage}\hfill
\begin{minipage}[t]{0.32\linewidth}\writespace{3}\end{minipage}\hfill
\begin{minipage}[t]{0.32\linewidth}\writespace{3}\end{minipage}

\section{마무리 성찰}

\begin{sheetbox}
\begin{minipage}[t]{0.48\linewidth}
{\footnotesize\color{inksoft} 일대일함수와 일대일대응의 차이를 한 문장으로}
\writespace{2}
\end{minipage}\hfill
\begin{minipage}[t]{0.48\linewidth}
{\footnotesize\color{inksoft} 감사 일기 / 유념 한 줄}
\writespace{2}
\end{minipage}
\end{sheetbox}

\end{document}
